\scanpage{9}
qui s'explicitent, en termes des $\varphi_i=f_i g_i$, et des $\alpha_i,\lambda_i$ correspondants, ainsi
\[
 \left\{\begin{aligned}
 \bar\alpha_k&=\alpha_k(X_k)\circ\mathrm{pr}_k\\
 \bar\lambda_{k,j,i}&=\lambda_{kji}(X_i)\circ\mathrm{pr}_i
 \end{aligned}\right. \tag*{(3.9)}
\]
\diagram{\((3.10)\quad\boxed{\begin{tikzcd}[ampersand replacement=\&,column sep=1.3em,row sep=.25em]
 \lambda_{kji}: \& \varphi_{ki}\arrow[r] \& \varphi_{kj}\circ\varphi_{ji}\\
 \& \shortparallel \& \shortparallel\\
 \& f_k g_i \& f_k\underbrace{g_j f_j}g_i\\
 \& \& \mathrm{id}\arrow[u]
 \end{tikzcd}}\)}
Ces $(\varphi_i,\lambda_{ijk})$ sont reliés par des conditions exprimant (1.5) (1.6), que nous n'explicitons pas. \inserted{Pour \deleted{\unsure{k}}$i=j=k$, $\lambda_{kji}$ n'est autre que $\lambda_i$.} Le cas le plus intéressant est celui où les \deleted{\unsure{$\varphi_i$|$\varphi$}} $\alpha_i$ sont des isom, (i.e. les $g_i$ \unsure{pl|plt} fidèles).
% slanted
\margin{notons que si deux des indices $i,j,k$ coïncident, alors $\lambda_{kji}$ est trivialement donné en termes des $\lambda_i$ --}

Dans ce cas, \deleted{il est} on peut supposer $\varphi_i=\mathrm{id}_{B_i}$, $\alpha_i=\mathrm{id}_{\mathrm{id}_{B_i}}$, $\lambda_i=\mathrm{id}$. On constate (!) que \inserted{dans} le cas où deux \inserted{consécutifs} des indices $i,j,k$ sont égaux, i.e. $i=j$ ou $j=k$, de sorte que l'un des facteurs de $\varphi_{kj}\varphi_{ji}$ est l'identité, on trouve \inserted{$\lambda_{kji}$} l'\uline{identité}.

Il reste donc à examiner les cas
\[
 \left\{\begin{aligned}
 \text{a)}\ &\deleted{\unsure{i\ne j,\ j\ne k,\ i\ne{}}}\quad i,j,k\ \text{tous différents}\\
 \text{b)}\ &i=k\ne j\qquad\mathrm{id}_{B_i}\longrightarrow\varphi_{ij}\circ\varphi_{ji}\qquad(i\ne j)
 \end{aligned}\right. \tag*{(3.11)}
\]
\deleted{Si \unsure{nous}}

\centerline{\rule{.72cm}{.3pt}\quad\rule{.72cm}{.3pt}\quad\rule{.72cm}{.3pt}\quad\rule{.9cm}{.3pt}}
% archive sheet number: 8
